% ===== §3 The First Cause: Encoding Mismatch =====
\section{The First Cause: Encoding Mismatch}
\label{sec:encoding}

\noindent
The etiology begins where culpability is least. Before power has accumulated, before any defence has been mounted or any advantage pressed, two subjects who wish each other nothing but good can still pass without meeting---and they can do so for a reason that lies entirely in the medium through which they reach one another, not in any defect of will. This first cause is the mismatch of the languages in which trust and care are encoded. It is the most intimate of the causes and the one that most resembles innocence, and for exactly that reason it is the one most often misread, by those living through it, as betrayal.

\subsection{Sincerity has no universal language}
\label{sub:encoding-nolang}

The series has already established, in its treatment of the non-binding vow, a thesis that the present section turns to a new use: that \emph{sincerity has no universal language}.\footnote{See \citet{paperxiii}, the chapter on translation: sincerity possesses no medium-independent form but must be encoded, and the encoding language that will succeed depends on the epistemology of its receiver. The present section runs that mechanism in reverse, asking what happens when the encoding fails.} A subject who means well must nonetheless \emph{express} the meaning well, and there is no expression that is meaning as such, prior to and independent of any code. Care must be put into a form---a word, a gift, an act, a silence---and every such form is legible only against a background grammar that assigns it sense. The trouble is that the grammar is not shared. What counts, for one subject, as the very signature of devotion---the reliable provision of material security, say, or the scrupulous keeping of small promises, or the offering of unsolicited help, or the granting of space---may register, against the other's grammar, as something between neutral and faintly insulting. The signal is sent in good faith and arrives as noise, or worse, as a signal of the opposite of what was meant.

\begin{claim}[The shadow of sincerity]
Between two subjects, sincerity is never transmitted directly; it casts a \emph{symbolic shadow}---the encoded, expressible form by which alone it can be conveyed---and it is the shadow, not the sincerity, that the other receives and must read. When the two grammars differ, the shadow cast by one subject's devotion falls, on the other's wall, in a shape the other cannot recognise as devotion. Estrangement of this first kind is not a failure of love but a failure of the shadow to be read---a mismatch of codes, prior to and independent of any fault.
\end{claim}

\subsection{The grammars of trust}
\label{sub:encoding-grammars}

The point is sharpened when the languages in question are specifically the languages of \emph{trust}---the grammars by which each subject decides what makes another worthy of being relied upon. The series has distinguished a family of such grammars, mutually intelligible only with effort and often not at all. One subject may trust on broadly \emph{evidential} grounds, extending reliance in proportion to a track record of verified performance, and reading the demand for such a record as simple prudence. Another may trust on grounds of \emph{character or virtue}, extending reliance to a person judged to be of a certain kind, and experiencing the demand for evidence as an insult to a bond that ought to be above audit. A third may trust \emph{inductively and relationally}, on the slow accumulation of shared history, so that what builds trust is neither proof nor judgment of character but sheer continued presence. A fourth may trust \emph{institutionally}, reposing confidence in the forms and guarantees---the vow witnessed, the commitment registered, the role formally assumed---and finding a trust that rests on none of these to be alarmingly unmoored.

These are not degrees of the same trust but different grammars of what trust \emph{is}, and a subject fluent in one is, in the others, not merely unpractised but apt to misread. To the evidential subject, the character-truster's refusal to keep score looks like carelessness; to the character-truster, the evidential subject's score-keeping looks like the absence of real faith. Each, acting from within an unimpeachable conception of what trusting well requires, presents to the other the very face of \emph{failing} to trust. Here the symptom of decoherence catalogued in \xref{sub:symptom-three} finds its first cause: a trust that is whole on each side can read, across the grammar-gap, as a trust that is draining away.

\subsection{Why this cause precedes responsibility}
\label{sub:encoding-precedes}

Two features of encoding mismatch must be fixed before the diagnosis proceeds, because both will be needed by the ethics and the praxis later.

The first is that this cause \emph{precedes responsibility}. There is, at this layer, no wrong yet done. Neither grammar is mistaken; neither subject has failed in care or in good faith; the misalignment is a property of the pairing, not of either party. This is of consequence for the normative argument (\xref{sec:asymmetry}), because it means the earliest and one of the commonest forms of estrangement is one to which the language of fault simply does not yet apply. To reach for blame here---to read the unread shadow as evidence of the other's coldness or one's own insufficiency---is not merely unkind but \emph{diagnostically false}: it imports a category, responsibility, that the situation does not contain, and in doing so converts a reparable mismatch into a manufactured grievance. A principal harm at this layer is done not by the mismatch but by its misreading.

The second is that encoding mismatch, though it precedes power, is the \emph{site} at which power will later take hold. For a grammar is never merely private. Which grammar of trust gets treated as the default---the ``normal'' or ``mature'' or ``reasonable'' way to trust, against which the other's grammar appears as deviation requiring explanation---is not decided by nature but settled, in any given dyad and any given culture, by forces that the later sections will name. When one subject's trust-grammar is tacitly installed as the standard and the other's marked as the exception, the difference of codes has already begun to harden into an asymmetry of standing. The mechanism of that hardening is not yet our subject; its material and historical sources are reserved for \xref{sec:histmat}. What matters here is only to mark the seam: the most innocent of the causes is also the surface along which the least innocent will eventually run.

\subsection{The corrective, and its limit}
\label{sub:encoding-corrective}

If the cause is mismatch, the corrective is, in principle, translation---and the series has already described its mature form. Between incommensurable grammars, trust can be rebuilt neither by one subject projecting their own code onto the other (which merely repeats the mismatch louder), nor by demanding the other adopt one's own (which is mismatch backed by force). It can be rebuilt only by \emph{mutual translation}, in which each labours to render their own sincerity into a shadow the other can read, while extending to the other's alien grammar what the series has called \emph{epistemic hospitality}---the presumption that the other's way of trusting is a way of trusting, and not a deficiency in it.\footnote{On the three paths---projection, multiple encoding, and mutual translation---and on epistemic hospitality as the corrective to both manipulation and a covert epistemological arrogance, see \citet{paperxiii}.}

But the corrective has a limit that it is the special business of \emph{this} paper, and not of its predecessor, to mark. Translation is a labour, and labour is performed under conditions; and where the conditions are set by an asymmetry of power---where one subject can afford to wait for the other to translate and the other cannot, or where one grammar is backed by the authority of a family, a custom, or a purse and the other is not---the ideal of mutual translation does not simply fail to obtain. It \emph{inverts}. What presents as translation becomes the quiet imposition of the stronger grammar under the name of mutual understanding; the weaker party, translating ever harder into a code not their own, mistakes their growing fluency in the other's language for the growth of intimacy, when it is in fact the early accomplishment of domination. The corrective to encoding mismatch is therefore available in full only between subjects whose power is not too unequal---which is to say, only when the conditions that the rest of the paper concerns itself with are already, to some degree, secured. With that limit, the first cause hands the inquiry to the second: for even where the grammars are matched and the will to translate is whole, two subjects can still fall out of phase, by a mechanism that lies not in the code but in the act of prediction itself.
